% !TeX program = pdflatex
\documentclass[11pt,a4paper,oneside]{book}

\usepackage[T1]{fontenc}
\usepackage[utf8]{inputenc}
\usepackage[polish]{babel}
\usepackage{lmodern}
\usepackage{amsmath,amssymb,amsthm,mathtools}
\usepackage{microtype}
\usepackage{xcolor}
\usepackage{geometry}
\usepackage{enumitem}
\usepackage{fancyhdr}
\usepackage[unicode]{hyperref}

\hypersetup{colorlinks=true,linkcolor=blue!45!black,urlcolor=blue!45!black}
\geometry{margin=2.5cm,headheight=15pt}
\setlength{\parindent}{0pt}
\setlength{\parskip}{4pt}
\setlength{\emergencystretch}{2em}
\setlist[enumerate]{topsep=4pt,itemsep=2pt}
\setlist[itemize]{topsep=4pt,itemsep=2pt}
\numberwithin{equation}{section}

\theoremstyle{plain}
\newtheorem{twierdzenie}{Twierdzenie}[chapter]
\newtheorem{lemat}[twierdzenie]{Lemat}
\newtheorem{wniosek}[twierdzenie]{Wniosek}
\theoremstyle{definition}
\newtheorem{definicja}[twierdzenie]{Definicja}
\newtheorem{przyklad}[twierdzenie]{Przykład}
\theoremstyle{remark}
\newtheorem{uwaga}[twierdzenie]{Uwaga}
\renewcommand{\proofname}{Dowód}

\pagestyle{fancy}
\fancyhf{}
\fancyhead[L]{\small Teoria Galois}
\fancyhead[R]{\small\thepage}
\renewcommand{\headrulewidth}{0.4pt}
\fancypagestyle{plain}{%
  \fancyhf{}
  \fancyfoot[C]{\thepage}
  \renewcommand{\headrulewidth}{0pt}
}
\newcommand{\Q}{\mathbb{Q}}
\newcommand{\R}{\mathbb{R}}
\newcommand{\C}{\mathbb{C}}
\newcommand{\Z}{\mathbb{Z}}
\newcommand{\F}{\mathbb{F}}
\newcommand{\K}{\mathbb{K}}
\newcommand{\Gal}{\operatorname{Gal}}
\newcommand{\Aut}{\operatorname{Aut}}
\newcommand{\Tr}{\operatorname{Tr}}
\newcommand{\Nm}{\operatorname{N}}
\newcommand{\ord}{\operatorname{ord}}
\newcommand{\charac}{\operatorname{char}}
\newcommand{\id}{\operatorname{id}}
\newcommand{\res}{\operatorname{res}}
\newcommand{\dd}{\,\mathrm{d}}
\title{Teoria Galois\\\large Rozszerzenia ciał, równania wielomianowe,\\ konstrukcje geometryczne i całki elementarne}
\author{}
\date{6 października 2026}
\begin{document}
\frontmatter
\maketitle
\tableofcontents
\mainmatter

\input{../rozdzialy/teoria_galois/00-ciala-i-rozszerzenia.tex}
\input{../rozdzialy/teoria_galois/01-ciala-rozkladu-i-grupy-galois.tex}
\input{../rozdzialy/teoria_galois/02-odpowiedniosc-galois.tex}
\input{../rozdzialy/teoria_galois/03-grupy-rozwiazalne-i-rozszerzenia-przez-pierwiastniki.tex}
\input{../rozdzialy/teoria_galois/04-dlaczego-nie-ma-ogolnego-wzoru-od-stopnia-piatego.tex}
\input{../rozdzialy/teoria_galois/05-konstrukcje-cyrklem-i-linijka.tex}
\input{../rozdzialy/teoria_galois/06-dalsze-przyklady-cial-i-zastosowan.tex}
\input{../rozdzialy/teoria_galois/07-ciala-rozniczkowe-i-twierdzenie-liouville-a.tex}
\input{../rozdzialy/teoria_galois/08-cztery-wazne-calki-nieelementarne.tex}
\input{../rozdzialy/teoria_galois/09-przyklady-z-rownan-rozniczkowych.tex}
\end{document}
